\scanpage{64}
\mspage{63}

on l'appelle le \uline{support complémentaire} du support donné $A$.

(vii) $\Sigma_{\unsure{\mathscr{F}|\mathcal{F}}}^\circ$ a un plus petit élément $\varnothing$ (à savoir la partie \deleted{vide} de $\mathscr{F}$\inserted{ réduite à $\varnothing_{\mathscr{F}}$*}), et un plus grand élément (à savoir $\mathscr{F}$), qui sont complémentaires l'un de l'autre.

% margin note written sideways; slanted upward
\margin{* Ceci provient du fait que $\varnothing_{\mathscr{F}}$ est l'unique élément de $\mathscr{F}$ qui soit $\parallel$ de tout élément de $\mathscr{F}$ (ou \unsure{encore} disjoint \unsure{d'elle-même|de lui-même}). On le notera, par abus de langage, $\varnothing$ (cf plus bas).\newline $\varnothing_\Sigma$}

(viii) \inserted{Plus gén.} Dans $\Sigma_{\mathscr{F}}^\circ$, il y a des Sup et des Inf quelconques, échangés par l'application $\unsure{\complement|\mathcal{C}}$. Le \deleted{\unsure{produit}} inf s'explicite par
\[
\underset{i}{\operatorname{Inf}}^{\Sigma_{\mathscr{F}}^\circ}A_i=\bigcap A_i
\]
et le Sup par
\[
\underset{i}{\operatorname{Sup}}^{\Sigma_{\mathscr{F}}^\circ}A_i=\overline{\bigcup A_i}\quad\Rightarrow\ (\text{ou }\operatorname{supp}(\bigcup A_i))
\]

\unsure{(iix)|(ix)} $\forall A\in\Sigma_{\mathscr{F}}^\circ$, on a
\[
A\cap\complement A=\varnothing_\Sigma
\]
\deleted{(En effet si $\unsure{F\in A}$ \unsure{de} $\varnothing_{\mathscr{F}}$}
(Exprimant le fait que le seul élément de $\mathscr{F}$ \deleted{\unsure{qui}} tel que $F\parallel F$ est $\varnothing_{\mathscr{F}}$, lequel est $\parallel$ de $\mathscr{F}$ tout entier\ldots)

($\unsure{x}$) \uline{Cor}
\[
A\subset B\ \Longrightarrow\ A\cap\complement B=\varnothing_\Sigma
\]

\uline{NB} Il n'est pas clair que l'implication inverse soit vraie. \inserted{$\Rightarrow$ \unsure{l'est}} Il n'est pas clair \deleted{que} (il est \deleted{parfois} faux) que le Sup\inserted{ de $\Sigma_{\mathscr{F}}^\circ$} \deleted{soit} s'interprète comme la réunion dans $\mathscr{F}$ -- la réunion\inserted{ (en tant que parties de $\mathscr{F}$)} de deux supports, n'est pas forcément un support : cela exprime le fait\inserted{ intuitif} qu'une ``figure''\inserted{ (dont le support est)} (\unsure{contenu}\deleted{\unsure{e}}) dans la réunion des supports de deux figures n'est pas forcément contenue dans\inserted{ (le support de)} l'une d'elles -- même dans le cas des multistrates, et $\unsure{=|\overline m}$ dans le cas des
